\chapter{Skills\\ทำให้ความเชี่ยวชาญมองเห็นได้}

คำว่า \enquote{มีทักษะการวิจัยและการสื่อสารที่ดี} บอกผู้อ่านได้น้อย เพราะไม่ระบุว่า Mentee ทำงานประเภทใด ด้วยเครื่องมือใด ในระดับใด และสร้างผลลัพธ์อะไร บทนี้ช่วยแปลงความสามารถที่ซ่อนอยู่ในประสบการณ์ให้เป็น Skills ที่ชัด พิสูจน์ได้ และพัฒนาต่อได้

\learningobjectives{
  \item ระบุ Skills จากพฤติกรรมและผลลัพธ์จริง
  \item แยกระดับความชำนาญออกจากการคุ้นเคย
  \item ออกแบบหลักฐานและแผนพัฒนา Skill ที่สำคัญ
}

\section{เขียน Skill ด้วยคำกริยา เครื่องมือ และบริบท}
Skill ที่ดีตอบว่า \textbf{ทำอะไร} \textbf{ด้วยอะไร} และ \textbf{ในเงื่อนไขใด} เช่น \enquote{ออกแบบแบบจำลอง computer vision สำหรับตรวจโรคพืชจากภาพภาคสนามที่มีแสงไม่สม่ำเสมอ} ชัดกว่า \enquote{มีทักษะ AI} และเปิดให้ตรวจสอบด้วยต้นแบบ ชุดข้อมูล หรือผลทดสอบ

\section{Skill Stack}
นวัตกรรมเกิดจากการประกอบ Skills มากกว่าทักษะเดี่ยว Mentee อาจมี technical skill ในการวิเคราะห์ข้อมูล, translational skill ในการออกแบบทดลองกับผู้ใช้, และ relational skill ในการประสานทีมหลายศาสตร์ การมอง Skill Stack ช่วยอธิบายคุณค่าที่เลียนแบบได้ยากและเห็นช่องว่างที่ขัดขวางโครงการ

\begin{table}[ht]
\centering
\begin{tabularx}{\textwidth}{P{.23\textwidth}YY}
\toprule
\textbf{กลุ่ม Skill} & \textbf{ตัวอย่าง} & \textbf{หลักฐาน}\\
\midrule
Technical & ออกแบบเซนเซอร์ วิเคราะห์ข้อมูล พัฒนาต้นแบบ & ผลทดสอบ โค้ด สิทธิบัตร\\
Translational & แปลโจทย์ผู้ใช้ ออกแบบ pilot วางเส้นทางกำกับ & บันทึกทดลอง การยอมรับใช้\\
Relational & นำทีม เจรจาความร่วมมือ สื่อสารข้ามศาสตร์ & ผลงานทีม หนังสือรับรอง\\
\bottomrule
\end{tabularx}
\end{table}

\section{ระดับความชำนาญที่ซื่อตรง}
การระบุระดับควรอิงพฤติกรรม: \textit{เรียนรู้} หมายถึงทำตามแนวทางได้; \textit{ปฏิบัติได้อย่างอิสระ} หมายถึงเลือกวิธีและแก้ปัญหาได้; \textit{ถ่ายทอดได้} หมายถึงอธิบายเหตุผลและพัฒนาผู้อื่นได้; \textit{สร้างแนวปฏิบัติใหม่} หมายถึงยกระดับมาตรฐานของสาขา ระดับเหล่านี้ช่วยให้ Mentee วางแผนพัฒนาโดยไม่ต้องกล่าวอ้างเกินจริง

\begin{examplebox}
แทนที่จะเขียนว่า \enquote{บริหารโครงการเก่ง} Mentee ระบุว่า \enquote{ออกแบบจังหวะการทำงานของทีมวิทยาศาสตร์และวิศวกรรม 8 คน ตั้งแต่กำหนดสมมติฐานถึง pilot ภาคสนาม ส่งมอบตรงเวลา 4 จาก 5 milestone} ข้อความนี้เผยทั้ง Skill บริบท และ Evidence
\end{examplebox}

\diagnostic{หากต้องส่งต่อหนึ่งส่วนของงานให้คนใหม่ เขาต้องเรียนรู้อะไรจาก Mentee ก่อนจึงจะทำได้ดี}

\begin{mentornote}
อย่าจำกัด Skills เฉพาะสิ่งที่ตลาดซื้อได้ทันที งานพื้นฐาน การสร้างทีม และความสามารถเชิงจริยธรรมก็มีคุณค่า เพียงต้องอธิบายกลไกที่ทำให้เกิดคุณค่านั้นอย่างตรงไปตรงมา
\end{mentornote}

\begin{keytakeaway}
Skills มีความหมายเมื่อเขียนเป็นความสามารถที่สังเกตได้ เชื่อมกับบริบท มีระดับความชำนาญที่ซื่อตรง และมีหลักฐานสนับสนุน
\end{keytakeaway}

\begin{reflection}
เลือกโครงการหนึ่งและเขียน Skills ที่ใช้จริง 8 ข้อ จัดกลุ่มเป็น technical, translational และ relational เลือก Skill หนึ่งข้อที่เป็นคอขวด แล้วกำหนดประสบการณ์ที่จะช่วยพัฒนาใน 90 วัน
\blankline\blankline
\end{reflection}

\chaptersummary{การทำให้ Skills มองเห็นได้ช่วยให้ Mentee เข้าใจทุนความสามารถของตน เลือกผู้ร่วมงาน และวางแผนพัฒนาอย่างมีหลักฐาน Mentor มีหน้าที่ช่วยแยกคำกล่าวกว้าง ๆ ออกจากความสามารถที่สังเกตและถ่ายทอดได้}
